\documentclass[10pt,a4paper]{amsart}
\usepackage[margin=19mm]{geometry}
\usepackage{amsmath,amssymb,mathtools}
\usepackage[scr=rsfs,cal=euler]{mathalfa}
\usepackage{xfrac}
\usepackage[unicode,hidelinks,bookmarks=false]{hyperref}
\newcommand{\ie}{i.e.\ }
\newcommand{\cf}{cf.\ }
\newcommand{\resp}{respectively\ }
\newcommand{\Subsection}{\subsection}
\providecommand{\nobreakdash}{\nobreak\hbox{-}}
\setlength{\emergencystretch}{2em}
\multlinegap0pt
\usepackage{amssymb}
\usepackage{mathtools}
\usepackage{xcolor}
\usepackage{float}
\newtheorem{proposition}{Proposition}
\title{A Superdiffusive Local Limit Theorem for the Elephant Random Walk\\ and the Breakdown of Log-Concavity}
\author{H\'elio Trinas, Glauco Valle}
\date{Source: arXiv:2609.14753v1 (CC BY 4.0)}
\begin{document}
\maketitle

\noindent\emph{Source and licence.} A Superdiffusive Local Limit Theorem for the Elephant Random Walk\\ and the Breakdown of Log-Concavity. Version arXiv:2609.14753v1 (CC BY 4.0), distributed by arXiv under the Creative Commons Attribution 4.0 International licence, http://creativecommons.org/licenses/by/4.0/, as recorded in the arXiv licence metadata for this exact version and shown on the abstract page https://arxiv.org/abs/2609.14753v1. Full licence notice: \texttt{licenses/arxiv-2609.14753-CC-BY-4.0.txt}.\par\smallskip

\noindent\textit{Editorial scope.} Counted unit: the article's proposition prop:lattice (source0.tex lines 268--276). The article's complete original proof of it is delivered with it (source0.tex lines 278--358, one \texttt{proof} environment of 81 lines). No mathematics is written, repaired or re-derived: the statement and the proof are the article's own lines, in the article's own notation, and no retained line is substituted or re-flowed.  Retained uncounted as the source-local context the proof needs: lines 248--254; lines 256--264.  Nothing was omitted from any retained span, so the package carries no omission record.  No retained line contains a citation, so the wrapper carries no bibliography and no bibliography entry.  No retained line contains a figure environment, an image-inclusion command or drawing code, so no image is delivered and the package declares no asset dependency.  Editorial formatting only, in addition: an AMS \texttt{amsart} preamble that restates the article's own declarations and macros used by the retained lines, namely the environments \texttt{proposition} on separate counters, together with the standard packages those lines need.  The retained environments print in this document as Proposition 1 in order of appearance, whereas the article prints the same items with the numbering of its own full document.  The complete native source of the article as distributed in the arXiv e-print for this exact version is bundled unchanged as \texttt{sources/210352/source0.tex}.
\begin{equation}\label{eq:massrec}
u_{n+1,k}
=
A_{n,k}u_{n,k}
+
B_{n,k}u_{n,k-1},
\end{equation}

\begin{equation}\label{eq:coefficients}
A_{n,k}
=
\frac{1+a}{2}-\frac{ak}{n},
\qquad
B_{n,k}
=
\frac{1-a}{2}+\frac{a(k-1)}{n}.
\end{equation}

\begin{proposition}\label{prop:lattice}
For each fixed $a\in(1/2,1)$ there are finite positive constants $C_0,C_1$ such that,
for every $q\in[0,1]$ and $n\geq1$,
\begin{equation}\label{eq:lattice-bounds}
 \max_k u_{n,k}\leq C_0n^{-a},
 \qquad
 \max_k|u_{n,k}-u_{n,k-1}|\leq C_1n^{-2a}.
\end{equation}
\end{proposition}

\begin{proof}
Set $M_n=\max_k u_{n,k}$. For $1\leq k\leq n$, both coefficients in
\eqref{eq:massrec} are nonnegative and by \eqref{eq:coefficients}
\[
 A_{n,k}+B_{n,k}=1-\frac an.
\]
Also recall that at $k=0$, $u_{n,-1}=0$ and $A_{n,0}=p=(1+a)/2$. At $k=n+1$,
$u_{n,n+1}=0$ and $B_{n,n+1}=p$. Since there exists $N_0 = N_0(a)$ such that $p \le 1- a/n$ for $n\ge N_0$, we also have 
\begin{equation}\label{eq:mass-contraction}
 M_{n+1}\leq\left(1-\frac an\right)M_n, \ \forall \, n\ge N_0.
\end{equation}

Put $d_{n,k}=u_{n,k}-u_{n,k-1}$. Subtracting the recurrence at $k-1$ from
the recurrence at $k$ gives
\begin{align*}
 d_{n+1,k}
 &=A_{n,k}u_{n,k}+(B_{n,k}-A_{n,k-1})u_{n,k-1}
   -B_{n,k-1}u_{n,k-2}.
\end{align*}
The identities
\[
 A_{n,k-1}=A_{n,k}+\frac an,
 \qquad B_{n,k}=B_{n,k-1}+\frac an,
\]
imply
\[
 B_{n,k}-A_{n,k-1}=B_{n,k-1}-A_{n,k}.
\]
Substitution and regrouping therefore yield the exact recurrence
\begin{equation}\label{eq:difference-rec}
 \begin{aligned}
 d_{n+1,k}
 &=A_{n,k}(u_{n,k}-u_{n,k-1})
   +B_{n,k-1}(u_{n,k-1}-u_{n,k-2})\\
 &=A_{n,k}d_{n,k}+B_{n,k-1}d_{n,k-1}.
 \end{aligned}
\end{equation}
The boundary cases of \eqref{eq:difference-rec} follow directly from the
definition of $u_{n,k}$. For
$1\leq k\leq n+1$,
\[
 A_{n,k}+B_{n,k-1}=1-\frac{2a}{n}.
\]
The smallest active coefficient in this range is
$(1-a)/2-a/n$, which, enlarging $N_0 = N_0(a)$ if necessary, is nonnegative for all $n \ge N_0$. At $k=0$ in \eqref{eq:difference-rec}, only
$A_{n,0}d_{n,0}=p\,d_{n,0}$ remains; at
$k=n+2$, only
$B_{n,n+1}d_{n,n+1}=p\,d_{n,n+1}$ remains. Setting $D_n=\max_k|d_{n,k}|$, we have
\begin{equation}\label{eq:difference-contraction}
 D_{n+1}\leq\left(1-\frac{2a}{n}\right)D_n , \ \forall \, n\ge N_0(a).
\end{equation}

If $y_{n+1}\leq(1-c/n)y_n$ for $n\geq N_0 >c$, then
\[
 y_n\leq y_{N_0}\prod_{j=N_0}^{n-1}\left(1-\frac cj\right).
\]
Since $\log(1-x)\leq-x$ for $0<x<1$,
\begin{align*}
 \log\prod_{j=N_0}^{n-1}\left(1-\frac cj\right)
 &\leq-c\sum_{j=N_0}^{n-1}\frac1j
 \leq-c\int_{N_0}^n\frac{dx}{x}
 =-c\log\frac nN_0.
\end{align*}
Exponentiating gives
\begin{equation}\label{eq:product-bound}
 \prod_{j=N_0}^{n-1}\left(1-\frac cj\right)\leq\left(\frac{N_0}{n}\right)^c.
\end{equation}

From \eqref{eq:product-bound}, $c=a$ in
\eqref{eq:mass-contraction} and
$c = 2a$ in \eqref{eq:difference-contraction}, 
% The index $N_0$ depends only on $a$.
% Moreover, $0\leq u_{n,k}\leq1$ implies $M_{N_0}\leq1$ and
% $D_{N_0}\leq1$ for every $q\in[0,1]$. Hence
\[
 M_n\leq M_{N_0} \frac{N_0^a}{n^a} \le \frac{N_0^a}{n^a}, \qquad
 D_n\leq D_{N_0} \frac{N_0^{2a}}{n^{2a}} \le \frac{N_0^{2a}}{n^{2a}},
 \qquad n\geq N_0.
\]
Thus the statement hold with $C_0 = N_0^a$ and $C_1 = N_0^{2a}$ which depend on $a$ but not on $q$.
\end{proof}

\end{document}
